\documentclass[11pt]{amsart}
\usepackage[T1]{fontenc}
\usepackage{lmodern,amsmath,amssymb,amsthm,microtype}
\usepackage[margin=1.15in]{geometry}
\usepackage[colorlinks=true,linkcolor=blue,citecolor=blue,urlcolor=blue]{hyperref}
% Repository locator for the public (non-anonymous) build of
% non_mf_groups_exist.tex.  Kept in a separate file so that the anonymous
% build has an anonymous source archive and not merely an anonymous PDF:
% exclude this file from a double-blind submission and the manuscript still
% compiles, printing "supplied with this submission as supplementary
% material" in place of the link.
%
% Loaded only when \anonymousfalse is set in the preamble.
\newcommand{\coderepository}{https://github.com/SauersML/group-approximation}
\newcommand{\coderepositoryname}{github.com/SauersML/group-approximation}

\hypersetup{pdftitle={An infinite simple Kazhdan sofic group},pdfauthor={\paperauthorname}}
\newtheorem{theorem}{Theorem}
\newtheorem{lemma}[theorem]{Lemma}
\newtheorem{corollary}[theorem]{Corollary}
\newcommand{\F}{\mathbb F}
\newcommand{\Z}{\mathbb Z}
\newcommand{\LC}{\operatorname{LC}}
\newcommand{\EL}{\operatorname{EL}}
\newcommand{\GL}{\operatorname{GL}}
\newcommand{\doi}[1]{\href{https://doi.org/#1}{\nolinkurl{doi:#1}}}
\raggedbottom
\emergencystretch=1.5em

\begin{document}
\title{An infinite simple Kazhdan sofic group}
\author{\paperauthorname}
\date{}
\subjclass[2020]{Primary 20E32; Secondary 20E26, 20F10, 22D10, 37B10, 16S35}

\begin{abstract}
Let $R$ be a finitely generated simple ring with local annihilation:
for every finite $F\subseteq R$, the elements $s$ with $tFs=0$ for some
nonzero $t$ span $R$. Then for $n\ge3$ the group $\EL_n(R)$ modulo its
centre is infinite, simple, and has property~\textup{(T)}. It is locally
embeddable into finite groups when $R$ embeds in an ultraproduct of
finite rings and has finite centre. Crossed products of infinite minimal
subshifts by $\Z$ over finite fields satisfy these hypotheses. So there
are infinite simple Kazhdan groups that are sofic and hyperlinear, which
answers the question of Brown and Ozawa and Pestov's sofic version of
it. Every Turing degree occurs as the word-problem degree of one of
these groups.
\end{abstract}
\maketitle

Can an infinite simple Kazhdan group be hyperlinear? Brown asked this
in its von Neumann algebra form in 2001~\cite[\S11, Question~7]{Brown},
and Ozawa in the hyperlinear form in 2003~\cite[p.~527]{Ozawa}.
Pestov's Open question~9.1 adds the sofic version~\cite{Pestov}.
The groups below answer all three forms positively.

A ring $R$ has \emph{local annihilation} if for every finite
$F\subseteq R$ the elements $s\in R$ with $tFs=0$ for some nonzero
$t\in R$ additively span $R$. Write $Z_n(R)$ for the centre of
$\EL_n(R)$.

\begin{theorem}\label{thm:ring}
Let $R$ be a finitely generated simple ring with local annihilation, and
let $n\ge3$. Then every normal subgroup of $\EL_n(R)$ is central or all
of $\EL_n(R)$, and $\EL_n(R)/Z_n(R)$ is an infinite, finitely generated,
simple group with property~\textup{(T)}. If $R$ is a unital subring of
an algebraic ultraproduct of finite rings, then $\EL_n(R)$ is locally
embeddable into finite groups \textup{(LEF)}, and so is
$\EL_n(R)/Z_n(R)$ when the centre of $R$ is finite.
\end{theorem}

\begin{theorem}\label{thm:main}
Let $X\subseteq A^{\Z}$ be an infinite minimal subshift over a finite
alphabet, with shift $T$, let $\F_q$ be a finite field, and let
$R_X=\LC(X,\F_q)\rtimes_T\Z$, where $\LC$ denotes locally constant
functions. Then $R_X$ satisfies all hypotheses of
Theorem~\ref{thm:ring}, and its centre is $\F_q$. So for every $n\ge3$
the group $\EL_n(R_X)/Z_n(R_X)$ is infinite, finitely generated, simple,
Kazhdan and LEF, and so it is sofic and hyperlinear. For $q=2$ the
centre $Z_n(R_X)$ is trivial, so
\[
  G_X=\EL_3\bigl(\LC(X,\F_2)\rtimes_T\Z\bigr)
\]
is itself such a group.
\end{theorem}

Property~\textup{(T)} follows from the theorem of Ershov and
Jaikin-Zapirain~\cite{EJZ}. The extraction of an elementary matrix in
Lemma~\ref{lem:normal} follows Stepanov~\cite[Lemma~4.3]{Stepanov}, and
local annihilation supplies the noncentral commutators it starts from.
The finite models of $R_X$ use periodic approximation, as in the proof
by Grigorchuk and Medynets that topological full groups of minimal
Cantor systems are LEF~\cite[Theorem~2.6]{GM}. For a hyperlinear
infinite simple group $G$ the factor $L(G)$ embeds in
$\mathcal R^\omega$~\cite[Proposition~7.1]{Ozawa}, so $G$ lies in the
unitary group of the separable McDuff factor
$L(G)\mathbin{\bar\otimes}\mathcal R$, which embeds in $\mathcal R^\omega$
as Brown's question asks. Thom constructed a finitely generated Kazhdan
LEF group that is not residually finite~\cite{Thom}, but his example is
not simple. A finitely presented LEF group is residually
finite~\cite{VershikGordon}, and an infinite simple group is not, so the
groups of Theorem~\ref{thm:main} are not finitely presented. We do not
know whether some finitely presented infinite simple Kazhdan group is
sofic.

\section{Normal subgroups and finite models}

For a unital ring $R$, $n\ge3$, $i\ne j$ and $r\in R$, put
$e_{ij}(r)=I_n+rE_{ij}$, where $E_{ij}$ is the matrix unit, and let
$\EL_n(R)$ be the subgroup of $\GL_n(R)$ generated by these matrices.
With $[g,h]=ghg^{-1}h^{-1}$,
\begin{equation}\label{eq:elementary}
 \begin{aligned}
 e_{ij}(r+s)&=e_{ij}(r)e_{ij}(s),\qquad i\ne j,\\
 [e_{ik}(r),e_{kj}(s)]&=e_{ij}(rs),\qquad i,j,k\text{ distinct}.
 \end{aligned}
\end{equation}
So if $R$ is generated as a ring by a finite set $B\ni1$, then $\EL_n(R)$
is generated by the matrices $e_{ij}(s)$ with $s\in B$. A central element
of $\EL_n(R)$ commutes with every $e_{ij}(1)$ and every $e_{ij}(r)$, so it
is $cI_n$ with $c\in Z(R)^\times$.

\begin{lemma}\label{lem:normal}
Let $R$ be a simple ring with local annihilation and $n\ge3$. Then every
normal subgroup of $\EL_n(R)$ is central or all of $\EL_n(R)$.
\end{lemma}

\begin{proof}
Let $N\trianglelefteq\EL_n(R)$ contain a noncentral $g$, and let $F$
consist of $1$ and the entries of $g$. By~\eqref{eq:elementary} the
coefficients $s$ with $[g,e_{ij}(s)]=1$ for all $i\ne j$ form an additive
subgroup, which is not $R$ because $g$ is not central. So there are $s$
and $t\ne0$ with $tFs=0$, and $i\ne j$, such that $\rho=[g,e_{ij}(s)]\ne1$.
Every entry of
\[
 \rho-I_n=\bigl(g\,sE_{ij}\,g^{-1}\bigr)e_{ij}(-s)-sE_{ij}
\]
is a sum of elements $fsr$ with $f\in F$ and $r\in R$, so $t\rho=tI_n$.

Put $Y=\rho^{-1}-I_n$, and choose an entry $Y_{qm}\ne0$ and an index
$l\notin\{m,q\}$. Since $R$ is simple, $Y_{qm}Rt\ne0$, so $z=Y_{qm}ct\ne0$
for some $c\in R$. The relation $t\rho=tI_n$ gives
$e_{ml}(ct)\rho=\rho+ctE_{ml}$, so
\[
 v=\rho^{-1}e_{ml}(ct)\,\rho\,e_{ml}(-ct)=I_n+Y\,ctE_{ml}\in N.
\]
Both $v$ and $v^{-1}$ differ from $I_n$ only in column $l$, $v_{ql}=z$,
and $(v^{-1})_{ll}$ is a unit. For $b\notin\{l,q\}$ we have
$e_{bq}(1)v=v\,e_{bq}(1)+zE_{bl}$, so
\[
 [e_{bq}(1),v]=e_{bl}\bigl(z\,(v^{-1})_{ll}\bigr)\ne I_n.
\]
By~\eqref{eq:elementary}, commutators with the matrices $e_{ij}(1)$ move
this element to every position. So the level
$J=\{r\in R:e_{pq}(r)\in N\text{ for all }p\ne q\}$ of $N$ is nonzero. It
is a two-sided ideal, since it is additive and, for $l\notin\{p,q\}$,
$e_{pq}(rs)=[e_{pl}(r),e_{lq}(s)]$ and $e_{pq}(sr)=[e_{pl}(s),e_{lq}(r)]$.
So $J=R$ and $N=\EL_n(R)$.
\end{proof}

A group is LEF~\cite{VershikGordon} if every finite subset $F$ has a
finite model, that is, an injective map into a finite group that
preserves all products staying in $F$. LEF groups are sofic, and sofic
groups are hyperlinear~\cite[Example~4.5 and Theorem~3.3]{Pestov}.

\begin{lemma}\label{lem:lef}
Let $R$ be a unital subring of an algebraic ultraproduct
$\prod_\omega Q_k$ of finite rings, and let $n\ge3$. Then $\EL_n(R)$ is
LEF, and so is $\EL_n(R)/Z$ for every finite central subgroup $Z$.
\end{lemma}

\begin{proof}
The group $\EL_n(R)$ is a subgroup of $\prod_\omega\GL_n(Q_k)$. Finitely
many products and inequalities among its elements hold in $\GL_n(Q_k)$
for $\omega$-almost every $k$, so one coordinate gives a finite model of
any finite set. For the quotient, let $F\subseteq\EL_n(R)$ be finite with
$1\in F=FZ$, and let $\psi$ be a finite model of $FF$ in a group $Q$.
Then $\psi(Z)$ is a subgroup that centralizes $\psi(F)$, and
$fZ\mapsto\psi(f)\psi(Z)$ is a finite model of $F/Z$ in
$\langle\psi(F)\rangle/\psi(Z)$.
\end{proof}

\begin{proof}[Proof of Theorem~\ref{thm:ring}]
By~\eqref{eq:elementary}, $\EL_n(R)$ is finitely generated, and it has
property~\textup{(T)} by~\cite[Theorem~1.1]{EJZ}, since $R$ is finitely
generated. Property~\textup{(T)} passes to quotients, and
Lemma~\ref{lem:normal} makes $\EL_n(R)/Z_n(R)$ simple. If $R$ were
finite, then $F=R$ would give $tRs=0$ with $t\ne0$, which forces $s=0$
in a simple ring. So $R$ is infinite, and $e_{12}$ embeds $R$ in
$\EL_n(R)/Z_n(R)$ because central elements are scalar. The last
statement is Lemma~\ref{lem:lef}, since
$Z_n(R)\subseteq Z(R)^\times I_n$.
\end{proof}

\section{Subshift algebras}

Let $X$, $T$, $\F_q$ and $R=R_X$ be as in Theorem~\ref{thm:main}. Write
$e_U$ for the indicator of a clopen set $U\subseteq X$. Our conventions
are $(Tx)_t=x_{t+1}$ and
\[
 R=\bigoplus_{j\in\Z}\LC(X,\F_q)\,u^j,\qquad ufu^{-1}=f\circ T^{-1},
\]
so every element of $R$ is uniquely a finite sum $\sum_jf_ju^j$, and
$u^je_Uu^{-j}=e_{T^jU}$. Minimality and infiniteness imply that $T$ has
no periodic points.

\subsection*{Generation, simplicity, centre}
Products of the translates $u^je_au^{-j}$ of the letter indicators
$e_a=e_{\{x:x_0=a\}}$ are the cylinder indicators, which span
$\LC(X,\F_q)$. So $R$ is generated by $u^{\pm1}$, the $e_a$ and a
generator of $\F_q$.

A nonzero ideal $I$ contains an element $r=\sum_jf_ju^j$ with $f_0\ne0$,
after multiplication by a power of $u$. Choose a nonempty clopen $U$ on
which $f_0$ is a nonzero constant $c$ and with $U\cap T^jU=\varnothing$
for the nonzero exponents $j$ of $r$. Then $e_Ure_U=ce_U$, so $I$ contains
$e_U$ and its translates $e_{T^jU}$. By minimality and compactness
finitely many of them cover $X$, so $1=1-\prod_i(1-e_{T^{j_i}U})\in I$.

If $c=\sum_jc_ju^j$ is central, then
$e_Vc-ce_V=\sum_jc_j(e_V-e_{T^jV})u^j$ vanishes for every clopen $V$. A
clopen $V$ containing $x$ but not $T^{-j}x$ gives $c_j(x)=0$ for $j\ne0$.
Then $uc=cu$ makes $c_0$ invariant, so $c_0$ is constant by minimality,
and $Z(R)=\F_q$. In particular $Z_n(R)\subseteq\F_2^\times I_n=\{I_n\}$
when $q=2$.

\subsection*{Local annihilation}
Let $F\subseteq R$ be finite, with exponents at most $w$ in absolute
value. For $x\in X$ choose a nonempty clopen $W$ missing the points
$T^jx$ with $|j|\le w$, and then a clopen $V\ni x$ with
$W\cap T^jV=\varnothing$ for $|j|\le w$. Then
$e_Wfe_V=\sum_jf_je_{W\cap T^jV}u^j=0$ for $f\in F$, so $t=e_W$ gives
$tF(e_Vr)=0$ for all $r\in R$. By compactness such sets $V$ can be chosen
to partition $X$, and then $r=\sum_Ve_Vr$.

\subsection*{Finite models}
Fix $x\in X$ and $\ell\ge0$. By minimality every word of $X$ occurs in
$x$, and $x_{[0,2\ell)}$ recurs at arbitrarily large positions. Choose
such a position $m_\ell$ with all words of length $2\ell+1$ of $X$
occurring in $x_{[0,m_\ell+2\ell)}$. This word begins and ends with
$x_{[0,2\ell)}$, so it has period $m_\ell$, and the $m_\ell$-periodic
sequence $y_\ell$ that agrees with $x$ on $[0,m_\ell)$ has the same words
of length $2\ell+1$ as $X$. So $m_\ell$ is at least the number of these
words, which tends to infinity because $X$ is infinite.

On $\F_q^{\Z/m_\ell\Z}$ let $P\delta_t=\delta_{t+1}$. For $f$ depending only
on coordinates in $[-\ell,\ell]$ let $D_\ell(f)\delta_t=f(T^ty_\ell)\delta_t$,
where $f(T^ty_\ell)$ is the value of $f$ at any point of $X$ that agrees
with $T^ty_\ell$ on $[-\ell,\ell]$. Fix a nonprincipal ultrafilter $\omega$
on $\mathbb N$ and put
\[
 \varphi\Bigl(\sum_jf_ju^j\Bigr)=\Bigl(\sum_jD_\ell(f_j)P^j\Bigr)_\ell
 \in\prod_\omega M_{m_\ell}(\F_q),
\]
with $0$ in the finitely many coordinates where some $D_\ell(f_j)$ is
undefined. Since $PD_\ell(f)P^{-1}=D_\ell(f\circ T^{-1})$ once $\ell$
exceeds the windows involved, $\varphi$ is a unital ring homomorphism. If
$f_j\ne0$, then $D_\ell(f_j)\ne0$ for large $\ell$, and the terms with
$|j|<m_\ell/2$ have disjoint supports. So $\varphi$ is injective, and
Theorem~\ref{thm:ring} applies to $R_X$. This proves
Theorem~\ref{thm:main}.\hfill$\square$

\section{Word problems}

\begin{corollary}
The word problem of $G_X$ has the Turing degree of the language $L(X)$
of $X$. Every Turing degree occurs, so there are continuum many pairwise
nonisomorphic groups $G_X$.
\end{corollary}

\begin{proof}
Multiplying out a word in the generators gives a matrix with entries
$\sum_jf_ju^j$, each $f_j$ a table on words, using $uf=(f\circ T^{-1})u$.
The word is trivial if and only if the tables of its difference from
$I_3$ vanish on $L(X)$, so $L(X)$ computes the word problem. Conversely,
by~\eqref{eq:elementary} a word for $e_{12}(\prod_{t<n}u^{-t}e_{v_t}u^t)$
is computable from $v=v_0\cdots v_{n-1}$, and it is trivial if and only
if $v\notin L(X)$. For derived topological full groups, Grigorchuk and
Medynets proved that the word problem is decidable if and only if
$L(X)$ is recursive~\cite[Theorem~1.1(3)]{GMpres}.

For irrational $\alpha\in(0,1)$ let $X_\alpha$ be the infinite minimal
Sturmian subshift, the closure of the codings $c(\theta)$,
$\theta\in[0,1)$, with $c(\theta)_t=1$ if and only if
$\theta+t\alpha\bmod1\ge1-\alpha$~\cite{MorseHedlund}. Its words of
length $n$ are the constant values of $c(\theta)_{[0,n)}$ on the arcs
between the points $-j\alpha\bmod1$, $0\le j\le n$, so $\alpha$ computes
$L(X_\alpha)$. The word $c(\theta)_{[0,n)}$ has
$\lfloor\theta+n\alpha\rfloor$ ones, within $1$ of $n\alpha$, so
$L(X_\alpha)$ computes $\alpha$. Every set $S\subseteq\mathbb N$ has the
degree of the irrational number $[0;1+\chi_S(0),1+\chi_S(1),\dots]$, and
the degree of the word problem is an isomorphism invariant.
\end{proof}

\subsection*{Origin and authorship}
Claude (Anthropic) found the construction and original proofs under
the author's direction on September~12, 2026, and wrote the original
manuscript. Codex (OpenAI) shortened an earlier version, and Claude
revised this one.
The author is responsible for the final manuscript. Proof records are
in the \href{https://github.com/SauersML/group-approximation}{project repository}.

\begin{thebibliography}{99}

\bibitem{Brown}
N.~P. Brown, \emph{Tracial invariants, classification and
$\mathrm{II}_1$ factor representations of Popa algebras},
\href{https://arxiv.org/abs/math/0111286}{arXiv:math/0111286} (2001).

\bibitem{EJZ}
M.~Ershov and A.~Jaikin-Zapirain,
\emph{Property~\textup{(T)} for noncommutative universal lattices},
Invent. Math. \textbf{179} (2010), 303--347.
\doi{10.1007/s00222-009-0218-2}.

\bibitem{GM}
R.~Grigorchuk and K.~Medynets,
\emph{On algebraic properties of topological full groups},
Sb. Math. \textbf{205} (2014), 843--861.
\doi{10.1070/SM2014v205n06ABEH004400}.

\bibitem{GMpres}
R.~Grigorchuk and K.~Medynets,
\emph{Presentations of topological full groups by generators and
relations}, J. Algebra \textbf{500} (2018), 46--68.
\doi{10.1016/j.jalgebra.2016.10.027}.

\bibitem{MorseHedlund}
M.~Morse and G.~A. Hedlund,
\emph{Symbolic dynamics II. Sturmian trajectories},
Amer. J. Math. \textbf{62} (1940), 1--42.
\doi{10.2307/2371431}.

\bibitem{Ozawa}
N.~Ozawa, \emph{About the QWEP conjecture},
Internat. J. Math. \textbf{15} (2004), 501--530.
\doi{10.1142/S0129167X04002417}.

\bibitem{Pestov}
V.~G. Pestov, \emph{Hyperlinear and sofic groups: a brief guide},
Bull. Symbolic Logic \textbf{14} (2008), 449--480.
\doi{10.2178/bsl/1231081461}.

\bibitem{Stepanov}
A.~V. Stepanov, \emph{On the normal structure of the general linear
group over a ring}, Zap. Nauchn. Sem. POMI \textbf{236} (1997), 166--182;
English transl., J. Math. Sci. \textbf{95} (1999), 2146--2155.
\doi{10.1007/BF02169976}.

\bibitem{Thom}
A.~Thom, \emph{Examples of hyperlinear groups without factorization
property}, Groups Geom. Dyn. \textbf{4} (2010), 195--208.
\doi{10.4171/GGD/80}.

\bibitem{VershikGordon}
A.~M. Vershik and E.~I. Gordon,
\emph{Groups that are locally embeddable in the class of finite groups},
Algebra i Analiz \textbf{9} (1997), no.~1, 71--97; English transl.,
St. Petersburg Math. J. \textbf{9} (1998), no.~1, 49--67.

\end{thebibliography}
\end{document}
